% !TEX root = ../algebra_02.tex

\section{Действия групп}

\subsection{Определение и примеры}

\begin{Definition}{Действие группы на множество}{}
	Пусть $G$ — группа, $M$ — множество. \textbf{Левым действием} группы $G$ на множество $M$ (обозначается $G \curvearrowright M$) называется отображение $G \times M \to M$, заданное правилом $(g, m) \mapsto g \cdot m$ (или просто $gm$), удовлетворяющее двум свойствам:
	\begin{enumerate}
		\item $\forall g_1, g_2 \in G, \forall m \in M \colon (g_1 g_2)m = g_1(g_2 m)$.
		\item $\forall m \in M \colon em = m$, где $e$ — нейтральный элемент группы $G$.
	\end{enumerate}
\end{Definition}

\begin{Remark}{Правое действие}{}
	Аналогично определяется \textbf{правое действие} $M \times G \to M$, где $m \cdot g = mg$, со свойствами: $m(g_1 g_2) = (mg_1)g_2$ и $me = m$. Любое правое действие можно превратить в левое, определив $g \cdot m := m g^{-1}$.
\end{Remark}

\begin{Example}{Примеры действий групп}{}
	\begin{enumerate}
		\item Пусть $G = \mathbb{R}$, $M = \mathbb{C}$. Действие задается поворотом на угол: 
		\[ x \cdot z = z(\cos x + i \sin x) \]
		Проверим свойства:
		\begin{itemize}
			\item $x_1 \cdot (x_2 \cdot z) = z(\cos x_2 + i \sin x_2)(\cos x_1 + i \sin x_1) = z(\cos(x_1+x_2) + i \sin(x_1+x_2)) = (x_1+x_2) \cdot z$.
			\item $0 \cdot z = z(\cos 0 + i \sin 0) = z$.
		\end{itemize}
		
		\item $G = GL_n(K)$, $M = K^n$. Действие — умножение матрицы на вектор.
		
		\item $G \le S(M)$ (подгруппа симметрической группы). Действие задается естественным образом: $g \cdot m = g(m)$.
		
		\item $G = S_n$, $M = \{1, \dots, n\}^2$. Действие: $g \cdot (i, j) = (g(i), g(j))$. 
		
		\item Действие $G$ на себя левыми сдвигами: $M = G$. Действие $g \cdot m = gm$.
		
		\item Действие $G$ на себя сопряжениями: $M = G$. Действие $g \cdot m = gmg^{-1}$.
		Проверка свойств:
		\begin{itemize}
			\item $(g_1 g_2) \cdot m = (g_1 g_2)m(g_1 g_2)^{-1} = g_1(g_2 m g_2^{-1})g_1^{-1} = g_1 \cdot (g_2 \cdot m)$.
			\item $e \cdot m = e m e^{-1} = m$.
		\end{itemize}
	\end{enumerate}
\end{Example}

\subsection{Орбиты и стабилизаторы}

\begin{Definition}{Орбита}{}
	Пусть $G \curvearrowright M$. \textbf{Орбитой} элемента $m \in M$ называется множество:
	\[ Gm = \{gm \mid g \in G\} \]
\end{Definition}

\begin{Lemma}{О непересечении орбит}{}
	Для любых $m_1, m_2 \in M$ либо $Gm_1 = Gm_2$, либо $Gm_1 \cap Gm_2 = \emptyset$.
\end{Lemma}

\begin{proof}
	Предположим, что $Gm_1 \cap Gm_2 \neq \emptyset$. Тогда существует $y \in Gm_1 \cap Gm_2$, то есть найдутся $g_1, g_2 \in G$ такие, что $y = g_1 m_1 = g_2 m_2$. 
	Отсюда выразим $m_2$:
	\[ m_2 = g_2^{-1}(g_2 m_2) = g_2^{-1}(g_1 m_1) = (g_2^{-1}g_1)m_1 \]
	Тогда для любого $h \in G$ выполнено:
	\[ h m_2 = (h g_2^{-1}g_1)m_1 \in Gm_1 \]
	Следовательно, $Gm_2 \subset Gm_1$. Аналогично доказывается, что $Gm_1 \subset Gm_2$. Значит, $Gm_1 = Gm_2$.
\end{proof}

\begin{Remark}{Связь с классами эквивалентности}{}
	Орбиты являются классами эквивалентности на $M$ по отношению: $m_2 \sim m_1 \iff \exists g \in G \colon m_2 = gm_1$.
\end{Remark}

\begin{Definition}{Стабилизатор}{}
	Пусть $G \curvearrowright M$ и $m \in M$. \textbf{Стабилизатором} элемента $m$ называется множество:
	\[ St_m = \{g \in G \mid gm = m\} \]
	Очевидно, что $St_m \le G$ (является подгруппой), так как если $gm=m$, то умножая слева на $g^{-1}$, получаем $m = g^{-1}m$, то есть $g^{-1} \in St_m$.
\end{Definition}

\begin{Proposition}{Теорема об орбите и стабилизаторе}{}
	Если $G \curvearrowright M$, то мощность орбиты равна индексу подгруппы стабилизатора:
	\[ |Gm| = (G : St_m) \]
\end{Proposition}

\begin{proof}
	Построим отображение $\alpha \colon G/St_m \to Gm$, заданное правилом $\alpha(gSt_m) = gm$.
	\begin{itemize}
		\item \textbf{Корректность:} Если $g' = gh$ для некоторого $h \in St_m$, то $g'm = g(hm) = gm$. Отображение не зависит от выбора представителя смежного класса.
		\item \textbf{Сюръективность:} Очевидна по определению орбиты.
		\item \textbf{Инъективность:} Пусть $\alpha(g_1 St_m) = \alpha(g_2 St_m)$. Тогда $g_1 m = g_2 m \implies m = g_1^{-1}g_2 m \implies g_1^{-1}g_2 \in St_m \implies g_2 \in g_1 St_m \implies g_1 St_m = g_2 St_m$.
	\end{itemize}
	Построенная биекция доказывает утверждение.
\end{proof}

\begin{Corollary}{Следствие для конечных групп}{}
	Если группа $G$ конечна, то для любого $m \in M$ размер орбиты $|Gm|$ делит порядок группы $|G|$.
\end{Corollary}

\subsection{Действие как гомоморфизм}

\begin{Statement}{Эквивалентность действия и гомоморфизма}{}
	Задать действие группы $G$ на множестве $M$ — то же самое, что задать гомоморфизм из группы $G$ в симметрическую группу $S(M)$.
\end{Statement}

\begin{proof}
	Пусть $G \curvearrowright M$. Построим отображение $\alpha \colon G \to S(M)$, где $\alpha(g) = \alpha_g$, а $\alpha_g \colon M \to M$ действует как $m \mapsto gm$.
	Заметим, что $\alpha_g^{-1} = \alpha_{g^{-1}}$, поэтому $\alpha_g$ действительно является биекцией, то есть $\alpha_g \in S(M)$.
	Проверим свойство гомоморфизма:
	\[ \alpha_{g_1 g_2}(m) = (g_1 g_2)m = g_1(g_2 m) = \alpha_{g_1}(\alpha_{g_2}(m)) \]
	Следовательно, $\alpha_{g_1 g_2} = \alpha_{g_1} \circ \alpha_{g_2}$, и $\alpha$ — гомоморфизм.
	
	Обратно, если дан гомоморфизм $\alpha \colon G \to S(M)$, то действие определяется как $gm := \alpha_g(m)$, при этом нейтральный элемент переходит в $\alpha_e = \text{id}_M$.
\end{proof}

\begin{Theorem}{Теорема Кэли}{}
	Пусть $G$ — конечная группа порядка $n$. Тогда существует подгруппа $G' \le S_n$ такая, что $G \cong G'$.
\end{Theorem}

\begin{proof}
	Рассмотрим действие $G$ на себе левыми сдвигами. Оно определяет гомоморфизм $\alpha \colon G \to S(G)$, заданный как $g \mapsto (h \mapsto gh)$.
	Найдём ядро гомоморфизма:
	\[ \ker \alpha = \{g \in G \mid gh = h \quad \forall h \in G\} = \{e\} \]
	Так как ядро тривиально, $\alpha$ является инъективным гомоморфизмом. Зафиксировав биекцию между множеством $G$ и $\{1, 2, \dots, n\}$, мы индуцируем изоморфизм $\lambda \colon S(G) \xrightarrow{\sim} S_n$.
	Тогда композиция отображений $\lambda \circ \alpha$ задаёт изоморфизм группы $G$ на свой образ:
	\[ G \cong \operatorname{Im}(\lambda \circ \alpha) \le S_n \]
	Теорема доказана.
\end{proof}

\begin{Example}{Применение теоремы Кэли к четверной группе Клейна}{}
	Найдём изоморфную подгруппу $G' \le S_4$ для группы $G = (\mathbb{Z}/2\mathbb{Z}) \times (\mathbb{Z}/2\mathbb{Z})$.
	Элементы группы: $e=(0,0), a=(1,0), b=(0,1), c=(1,1)$.
	Занумеруем элементы: $1 \leftrightarrow e$, $2 \leftrightarrow a$, $3 \leftrightarrow b$, $4 \leftrightarrow c$.
	Действуя левым умножением каждого элемента на всю группу, получаем перестановки:
	\begin{itemize}
		\item Умножение на $e$ оставляет все на местах: $(1)(2)(3)(4) = \text{id}$.
		\item Умножение на $a$: $ae=a \to 2$, $aa=e \to 1$, $ab=c \to 4$, $ac=b \to 3$. Перестановка $(1\;2)(3\;4)$.
		\item Умножение на $b$: $be=b \to 3$, $ba=c \to 4$, $bb=e \to 1$, $bc=a \to 2$. Перестановка $(1\;3)(2\;4)$.
		\item Умножение на $c$: $ce=c \to 4$, $ca=b \to 3$, $cb=a \to 2$, $cc=e \to 1$. Перестановка $(1\;4)(2\;3)$.
	\end{itemize}
	Таким образом, $G \cong G' = \{\text{id}, (1\;2)(3\;4), (1\;3)(2\;4), (1\;4)(2\;3)\} \le S_4$.
\end{Example}

\subsection{Лемма Бёрнсайда}

\begin{Definition}{Неподвижные точки}{}
	Пусть $G \curvearrowright M$. Для элемента $g \in G$ множество его неподвижных точек обозначается как:
	\[ M^g := \{m \in M \mid gm = m\} \]
\end{Definition}

\begin{Lemma}{Лемма Бёрнсайда}{}
	Пусть $G$ и $M$ конечны. Тогда число орбит $N$ действия $G$ на $M$ равно среднему числу неподвижных точек:
	\[ N = \frac{1}{|G|} \sum_{g \in G} |M^g| \]
\end{Lemma}

\begin{proof}
	Рассмотрим множество пар $S = \{(g, m) \in G \times M \mid gm = m\}$. Посчитаем его мощность двумя способами.
	С одной стороны, суммируя по элементам группы:
	\[ |S| = \sum_{g \in G} |M^g| \]
	С другой стороны, суммируя по элементам множества:
	\[ |S| = \sum_{m \in M} |\{g \in G \mid gm = m\}| = \sum_{m \in M} |St_m| \]
	По теореме об орбите и стабилизаторе $|St_m| = \frac{|G|}{|Gm|}$. Подставим это в сумму:
	\[ \sum_{m \in M} \frac{|G|}{|Gm|} = |G| \sum_{m \in M} \frac{1}{|Gm|} \]
	Сгруппируем слагаемые по орбитам. Пусть $X = \{Gm \mid m \in M\}$ — множество всех орбит.
	\[ |G| \sum_{O \in X} \sum_{m \in O} \frac{1}{|O|} = |G| \sum_{O \in X} \left( |O| \cdot \frac{1}{|O|} \right) = |G| \sum_{O \in X} 1 = |G| \cdot |X| = |G| \cdot N \]
	Приравнивая два способа подсчета $|S|$, получаем:
	\[ \sum_{g \in G} |M^g| = |G| \cdot N \implies N = \frac{1}{|G|} \sum_{g \in G} |M^g| \]
\end{proof}

\begin{Example}{Раскраска ожерелья (группа $D_9$)}{}
	Рассмотрим задачу о количестве способов раскрасить правильный 9-угольник (ожерелье из 9 бусин) в 2 цвета с точностью до поворотов и отражений. Группа симметрий — диэдральная группа $D_9$, $|D_9| = 18$.
	Найдём мощности множеств неподвижных точек $|M^g|$ для каждого типа преобразований:
	\begin{itemize}
		\item $R_0$ (тождественное): 1 элемент. Все $2^9$ раскрасок остаются на месте. $|M^{R_0}| = 512$.
		\item $S_i$ (отражения): 9 элементов. Ось симметрии проходит через 1 вершину и середину противоположной стороны. Вершины разбиваются на 1 неподвижную и 4 пары симметричных. Для сохранения раскраски цвета в парах должны совпадать. Независимых выборов цвета: 5. $|M^{S_i}| = 2^5 = 32$.
		\item $R_i$ (повороты порядка 9): Это повороты на углы, кратные $2\pi/9$, взаимно простые с 9. Таких элементов $\varphi(9) = 6$. Орбита вершин одна, поэтому все вершины должны быть одного цвета. $|M^{R_i}| = 2^1 = 2$.
		\item $R_i$ (повороты порядка 3): Повороты на $120^\circ$ и $240^\circ$. Таких элементов 2. Вершины разбиваются на 3 независимых цикла по 3 вершины. $|M^{R_i}| = 2^3 = 8$.
	\end{itemize}
	Применим лемму Бёрнсайда:
	\[ N = \frac{1}{18} (1 \cdot 512 + 9 \cdot 32 + 6 \cdot 2 + 2 \cdot 8) = \frac{1}{18} (512 + 288 + 12 + 16) = \frac{828}{18} = 46 \]
	Всего существует 46 различных ожерелий.
\end{Example}

\subsection{Центр группы и $p$-группы}

\begin{Definition}{Центр группы}{}
	\textbf{Центром} группы $G$ называется множество элементов, коммутирующих со всеми элементами группы:
	\[ Z(G) = \{h \in G \mid \forall g \in G \colon gh = hg\} \]
	Очевидно, что $Z(G)$ является нормальной подгруппой $G$ ($Z(G) \triangleleft G$).
\end{Definition}

\begin{Theorem}{О центре конечной $p$-группы}{}
	Пусть $G$ — конечная $p$-группа, то есть $|G| = p^n$, где $p$ — простое число, $n \in \mathbb{N}$. Тогда центр группы нетривиален: $Z(G) \neq \{e\}$.
\end{Theorem}

\begin{proof}
	Рассмотрим действие $G$ на себе сопряжениями: $g \cdot m = gmg^{-1}$.
	Орбитами этого действия являются классы сопряженности. По следствию из теоремы об орбите и стабилизаторе, размер каждой орбиты $|Gm|$ делит $|G|$, поэтому $|Gm| = p^k$ для некоторого $0 \le k \le n$.
	Разбивая множество $G$ на орбиты, получаем уравнение классов:
	\[ |G| = \sum_{k=0}^{n} m_k p^k = p^n \]
	где $m_k$ — количество орбит размера $p^k$.
	Из этого равенства следует, что $p$ делит $m_0$ (количество орбит длины 1).
	Орбита элемента $m$ имеет длину 1 тогда и только тогда, когда $\forall g \in G \colon gmg^{-1} = m \iff gm = mg$, что означает $m \in Z(G)$.
	Следовательно, $m_0 = |Z(G)|$. Так как нейтральный элемент $e$ всегда лежит в $Z(G)$, имеем $|Z(G)| \ge 1$. Но так как $p$ делит $|Z(G)|$, мы получаем, что $|Z(G)| \ge p > 1$. Значит, $Z(G) \neq \{e\}$.
\end{proof}

\begin{Lemma}{Факторгруппа по центру}{}
	Если группа $G$ не является абелевой, то факторгруппа $G/Z(G)$ не может быть циклической.
\end{Lemma}

\begin{proof}
	Докажем от противного. Пусть $G/Z(G)$ — циклическая группа, порожденная смежным классом $g_0 Z(G)$:
	\[ G/Z(G) = \langle g_0 Z(G) \rangle \]
	Тогда любой элемент $g \in G$ представим в виде $g \in g_0^i Z(G)$ для некоторого целого $i$, то есть $g = g_0^i h$, где $h \in Z(G)$.
	Возьмем произвольные два элемента $g, g' \in G$. Они представимы как $g = g_0^i h$ и $g' = g_0^j h'$, где $h, h' \in Z(G)$.
	Рассмотрим их произведение:
	\[ gg' = (g_0^i h)(g_0^j h') = g_0^i g_0^j h h' = g_0^{i+j} h h' \]
	Мы смогли переставить элементы, так как $h$ и $h'$ лежат в центре и коммутируют со всем.
	Аналогично:
	\[ g'g = (g_0^j h')(g_0^i h) = g_0^j g_0^i h' h = g_0^{i+j} h h' \]
	Получили $gg' = g'g$ для любых $g, g' \in G$. Следовательно, $G$ — абелева группа. Противоречие с условием.
\end{proof}

\begin{Corollary}{Группы порядка $p^2$}{}
	Если $p$ — простое число, то любая группа $G$ порядка $p^2$ является абелевой.
\end{Corollary}

\begin{proof}
	Так как $G$ — $p$-группа, её центр нетривиален, следовательно, возможные порядки центра: $|Z(G)| = p$ или $|Z(G)| = p^2$.
	Предположим, $|Z(G)| = p$. Тогда индекс подгруппы $(G : Z(G)) = |G/Z(G)| = p$.
	Группа простого порядка всегда циклическая, значит $G/Z(G)$ циклическая. Но по предыдущей лемме, это влечет абелевость $G$.
	Если $G$ абелева, то она совпадает со своим центром, то есть $|Z(G)| = p^2$. Мы получили противоречие с допущением $|Z(G)| = p$.
	Следовательно, единственный возможный случай — это $|Z(G)| = p^2$, то есть $G = Z(G)$, и группа абелева.
\end{proof}